% ECONOMICS_PROBLEM: {"id": "G23-313", "title": "Platform credit and inclusion", "jel_code": "G23", "source_id": "OP-0328"}
% !TeX program = xelatex
\documentclass[11pt,a4paper]{article}
\usepackage[margin=23mm]{geometry}
\usepackage{fontspec}
\setmainfont{Latin Modern Roman}
\usepackage{amsmath,amssymb,mathtools}
\usepackage{xcolor,enumitem,booktabs,longtable,array,xurl,hyperref}
\definecolor{muted}{HTML}{53636C}
\hypersetup{colorlinks=true,urlcolor=blue,linkcolor=blue}
\setlength{\parindent}{0pt}
\setlength{\parskip}{5pt}
\newcommand{\R}{\mathbb R}
\newcommand{\E}{\mathbb E}
\newcommand{\N}{\mathbb N}
\newcommand{\ind}{\mathbf 1}
\newcommand{\Var}{\operatorname{Var}}
\newcommand{\Cov}{\operatorname{Cov}}
\newcommand{\supp}{\operatorname{supp}}
\DeclareMathOperator*{\argmin}{arg\,min}
\DeclareMathOperator*{\argmax}{arg\,max}
\newcommand{\entrymeta}[3]{{\sffamily\small\color{muted}\textbf{\texttt{#1}}\quad|\quad #2\par #3\par}}
\newcommand{\Problemref}[1]{\texttt{#1}}
\newcommand{\JELref}[2]{\texttt{#2} (JEL \texttt{#1})}
\begin{document}
\section*{Platform credit and inclusion}
\textbf{Problem ID:} \texttt{G23-313}\quad \textbf{JEL:} \texttt{G23}\par
% BEGIN ECONOMICS STATEMENT
\label{problem:OP-0328}
{\sffamily\small\color{muted}Original subject group: Banking, credit and financial stability\par}
\label{definitions:problem:30007610}
\textbf{Audit classification:} Explicit published open question. The BIS introduction explicitly identifies open policy trade-offs. The useful-credit definition and Pareto frontier are editorial; the initial attainable set was not by itself a frontier.
\subsection*{Definitions and precise research target}
Consider a two-sided lending platform linking borrowers \(i\) to lenders, using borrower data to screen risk. Let \(n\) measure the platform's network size, \(d\) its data access, and \(a\) a policy specifying interoperability, data portability, privacy restrictions and entry conditions. Equilibrium loan prices, approval probabilities and default losses depend on \((n,d,a)\), with platform entry and borrower participation endogenous. Define inclusion as access to loans that increase borrowers' expected utility net of repayment, not merely observed approval. Let \(W(a)\) combine borrower and lender surplus, real screening costs, privacy losses and systemic externalities. The task is to characterize the attainable outcome set and its Pareto frontier
\[\{(\text{useful credit access},\text{market power},\text{privacy loss},\text{default risk})(a):a\in A\}.\]
Here \(A\) is the stated feasible policy set and each component has a declared population and measurement rule. Identify when additional data improves screening and when it mainly strengthens incumbent advantage. Compare policies holding primitives fixed while allowing equilibrium adaptation. State the exclusions needed to identify data and network effects separately; selection into platforms can otherwise mimic improved lending. Determine the conditions under which portability raises inclusion without worsening risk or privacy outcomes.

\textbf{Additional formal definitions and scope (editorial).} Let \(V_i^{\rm loan}(a)\) and \(V_i^{\rm outside}(a)\) be borrower expected lifetime utilities under a declared preference, information and repayment/default model. Useful access is \(\Pr(\text{loan available},\,V_i^{\rm loan}(a)>V_i^{\rm outside}(a))\) for a fixed eligible population; behavioral borrowing and normative welfare can differ and require separate conventions. Define network size by active participants and data access by a specified vector of permissible features rather than one uninterpreted scalar. Measure market power using a structural markup, elasticity or concentration statistic chosen beforehand; concentration alone is not market power. The displayed attainable set is not yet a welfare frontier. A Pareto frontier retains points for which no feasible policy weakly improves access and weakly lowers each cost/risk component with at least one strict improvement. Components require consistent monetary or otherwise declared units.
\subsection*{Variants and assumption changes}
\begin{enumerate}
\item \textbf{Screening with competitive lenders (editorial).} Fix participation and network size and vary data access under a known borrower utility model.
\item \textbf{Endogenous platform scale (editorial).} Allow entry, network effects and portability, and compare Pareto frontiers for useful access, privacy and default risk under different conduct assumptions.
\end{enumerate}
\subsection*{References and source audit}
\begin{enumerate}
\item Primary BIS WP 986 cover verifies Croxson, Frost, Gambacorta and Valletti, January 2022; no DOI assigned. Locator: Introduction, printed p.2 (PDF p.4), second paragraph; Conclusion, p.26. Support: Platform policy trade-offs. Full PDF accessible. URL: \url{https://www.bis.org/publications/working-paper-986-platform-based-business-models-and-financial-inclusion.pdf}.
\end{enumerate}
\begin{enumerate}\raggedright
\item\hypertarget{definitions:source:30007610:1}{} Karen Croxson; Jon Frost; Leonardo Gambacorta; Tommaso Valletti. 2022. \emph{Platform-based business models and financial inclusion}. Introduction, printed p.2 (PDF p.4), second paragraph; Conclusion, p.26.
\url{https://www.bis.org/publications/working-paper-986-platform-based-business-models-and-financial-inclusion.pdf}.
\end{enumerate}

\subsection*{Common conventions: Source mathematical conventions}

These conventions apply to each entry unless explicitly replaced.

\textbf{Models and identification.} A primitive model \(m\) specifies all probability laws, feasible choices, information, equilibrium and measurement rules used by its target. The admissible class \(\mathcal M\) is fixed independently of outcomes. If \(P_O(m)\) is its observable probability law and \(\tau(m)\) the target, its identified set at observable law \(P\) is
\[
 \mathcal I_\tau(P)=\{\tau(m):m\in\mathcal M,\ P_O(m)=P\}.
\]
Point identification means that this set is a singleton. An empty set signals incompatibility of data and assumptions. Coordinate bounds need not describe the joint set. Bounds are sharp only if every retained target is attainable or arbitrarily approachable by compatible models. Finite bounds require explicit restrictions on unobserved quantities; unrestricted confounding, hidden wealth, extrapolation or arbitrary equilibrium selection can yield uninformative sets.

\textbf{Causal interventions.} A potential outcome \(Y_i(p)\) is the outcome under a complete policy regime \(p\), including timing, financing, information and market spillovers. The target population and units are fixed before treatment. With interference, use the complete assignment vector or a justified exposure mapping; individual no-interference notation alone cannot identify an economy-wide intervention. A difference of mean potential outcomes does not require the cross-world joint distribution. Conditioning on post-treatment migration, survival, enrollment or employment changes the estimand and needs separate justification.

\textbf{Statistical statements.} An estimator, confidence set, forecast or test specifies a sampling law, model class and loss/coverage criterion. Uniform \((1-\alpha)\)-coverage means \(\inf_{m\in\mathcal M}\Pr_m\{\tau(m)\in C_n\}\geq1-\alpha\), or an explicitly asymptotic version. The whole parameter space trivially has coverage, so informativeness requires a stated width, risk or power objective. A forecasting improvement is relative to a named benchmark, information set and evaluation population.

\textbf{Equilibrium and optimization.} An equilibrium comprises feasible plans, optimality, beliefs where needed, budget constraints and all market/resource clearing. The primitives or a stated benchmark define each correspondence \(\mathcal E(p;m)\). Nonemptiness is assumed or proved, never inferred just from compact policy sets. A selected-equilibrium target uses a declared selection rule; a robust target quantifies over all permitted equilibria. Use suprema or infima unless attainment is established. An \(\varepsilon\)-optimal policy is feasible and within \(\varepsilon\) of the finite optimum value. Report an empty feasible set separately.

\textbf{Welfare and accounting.} Normative utility, interpersonal normalization, social weights, numeraire, discounting and the use of public receipts are primitives. Behavioral utility need not coincide with the welfare criterion. Household welfare includes ownership income and rents once; adding profits again to compensating variation generally double-counts them. A monetary equivalent is defined by an explicit transfer and price convention, with monotonicity and range conditions. Average effects, mechanism contributions, transfers and welfare losses are different targets. Infinite sums require convergence and dynamic feasible plans require terminal or no-Ponzi conditions.

\textbf{Source versions.} A working-paper date, revision date and journal publication year are distinguished. A DOI for a journal version is not automatically the DOI of an earlier manuscript. Locators refer to the stated version and printed pages where specified. A metadata-only check does not verify a claim about a full-text conclusion. Every entry's source subsection narrows the provenance claim accordingly.
% END ECONOMICS STATEMENT
\end{document}
